\scanpage{207}
% archive sheet number: 206
Voici la généralisation de Tate.

$X$ variété non singulière \unsure{alg$/k$ alg clos}, $K=k(X)$

$A:\unsure{A(V)|A(X)}$ var. d'Albanese, \inserted{hom $\unsure{\varphi|\phi}:X\longrightarrow A\deleted{\unsure{(X)|(V)}}$}, $\deleted{\unsure{P|\mathcal P}}$ $\widehat A$ variété de Picard.

Soient $\unsure{N|\mathcal N}$ le noyau d'Albanese : $N=\operatorname{Ker}(\unsure{\mathcal Z_0(X)|Z_0(X)}\longrightarrow A(X))$

\deleted{\unsure{On veut}}

$\unsure{\mathcal I|\mathcal J|I}$ le groupe des \unsure{données|classes} de répartitions
\[
\vec f=(f_x)_{x\in X}
\]
telles qu'il existe un diviseur $D$ sur $X$, avec $(f_x)_x=[D]_x$ pour tt $x$. On a une application
\[
C=\mathcal I/\deleted{\unsure{U}}K^*\longrightarrow\unsure{\mathcal P(X)|P(X)}
\]
(classes de diviseurs), soit $\deleted{\unsure{\mathcal I_0}}C_0$ l'\unsure{image} inverse de $\widehat A$; \unsure{donc on a} un hom surjectif
\[
C_0\longrightarrow\widehat A
\]
$\mathcal I_0$ est l'image réciproque de $C_0$ dans $\mathcal I$.
On veut définir un accouplement
\[
N\times C_0\longrightarrow k^*
\]
et pour ceci, on définit un accouplement
\[
N\times\mathcal I_0\longrightarrow k^*
\]
qui est trivial sur $N\times K^*$. Soit
\[
\unsure{z|\zeta}=\sum a_i\unsure{\tilde x_i|x_i|\xi_i}\in N,\qquad\deleted{\unsure{\deg z=0|\deg=0}}
\]
\inserted{$\xi\in C_0$,} $\vec f=(f_x)_{x\in X}\in\mathcal I_0$ un représentant de $\xi$ \inserted{tel que $\operatorname{div}(\vec f)$ ne rencontre pas $z$},
$D$ le diviseur défini par $\vec f$, $\deleted{\unsure{\xi}}\dot D$ son \unsure{image} dans $\widehat A$, et $z'\in\mathcal Z_0(\widehat A)$ \deleted{\unsure{avec}}, $\deg z'=0$, tel que $\dot D=S(z')$.
Soient $\Theta$ un diviseur canonique sur $\deleted{X}A\times\widehat A$ en \unsure{bonne position} pour $\varphi(z)\times\deleted{\unsure{\xi|z}}z'$, et considérons
